\section{Model}\label{sec:model}

The economy is stationary. Time is continuous, but within a topic the relevant
literature is a flow, so we describe a researcher's choices per unit of time.

\subsection{Researchers, topics and time}

There is a continuum of topics. Topics that are open are ex ante identical. A
mass $N$ of researchers each has one unit of time per period, split into a
share $x$ spent writing and a share $r=1-x$ spent reading. Writing full time
yields $a$ papers (projects) per period. Topic $j$ hosts $n_j$ researchers who
write, and receives a flow of
\begin{equation}\label{eq:flow}
P_j = a\,x\,n_j
\end{equation}
papers per period. $P_j$ is the topic's \emph{density}. We call the ratio of
writing speed to reading speed, $a/b$, the \emph{pace} of science.

\subsection{Reading teaches}

A researcher who reads full time reads $b$ papers per period, so a reading
share $r$ yields $R=b\,r$ papers drawn from her topic's flow. What she reads
raises the quality of her own work by the factor $Q(R,P)$.

\begin{assumption}[Learning]\label{ass:Q}
$Q:\R_+^2\to[1,\infty)$ is continuous, twice continuously differentiable on
$\R_{++}^2$, nondecreasing in both arguments, with $Q(0,P)=Q(R,0)=1$ and
$\bar Q(R)\equiv\sup_P Q(R,P)<\infty$ for every $R$. Moreover $\ln Q$ is
jointly concave and supermodular on $\R_+^2$.
\end{assumption}

Nothing is learned without reading or without a literature to read, and what
one can learn is bounded by what one reads. Concavity of $\ln Q$ rules out
increasing returns to reading. Supermodularity says that reading is worth
more in a topic with more to learn from. Our running example is
\begin{equation}\label{eq:Qexample}
Q(R,P) = \big(1+\beta\,h(R,P)\big)^{\gamma},\qquad h(R,P)=\frac{RP}{R+P},\qquad \beta>0,\;\gamma\in(0,1].
\end{equation}
The harmonic aggregator $h$ is close to $R$ when the reader samples a small
part of a rich literature and close to $P$ when she reads a thin literature
exhaustively. Because $h$ is concave and homogeneous of degree one, the
example satisfies \cref{ass:Q} (\cref{app:proofs}).

\subsection{Reading protects}

Each paper in the flow \emph{covers} a given project independently with a small
probability. The number of covering papers is therefore
Poisson with mean $\theta P$. A covering paper is caught only if it is read. With reading
coverage $c=R/P$ it is caught with probability $\varphi(c)$.

\begin{assumption}[Detection]\label{ass:phi}
$\varphi:\R_+\to[0,1]$ is continuous and twice differentiable on $(0,\infty)$,
with $\varphi(0)=0$, $\varphi'>0$ and $\varphi''<0$.
\end{assumption}

Concavity says that the first papers one reads are the most informative about
overlap. Random reading with replacement gives $\varphi(c)=1-e^{-c}$. Our
example $\varphi(c)=1-e^{-\sigma c}$ indexes \emph{search precision} by
$\sigma$: a better search engine finds covering papers sooner.

\paragraph{Credit.} Overlap erodes credit multiplicatively. Each covering paper
that goes unread turns part of the project into a duplicate and removes a
share $\lambda$ of its credit. When a covering paper is caught, the
researcher chooses between two options. She can \emph{pivot} to a fresh variant at a time cost
that leaves a share $1-\kappa_p$. Or she can \emph{proceed} as a follow-on
that cites the earlier paper and keeps a share $\delta$ drawn from a
distribution $G$. She proceeds if $\delta>1-\kappa_p$. The expected share kept
after a detection is $1-\ell\equiv\E\max\{\delta,1-\kappa_p\}$.

\begin{assumption}[Credit]\label{ass:credit}
$0\le \ell<\lambda\le 1$ and $\theta>0$.
\end{assumption}

Detection is valuable ($\ell<\lambda$), and duplicates need not earn zero
($\lambda<1$ is allowed). The arXiv evidence (Fact~\ref{fact:credit}) requires
both. Credit here is private. \Cref{sec:catchup} separates it from social
value, since part of the credit lost to overlap may be a transfer to the
earlier paper. Let $\Delta\equiv\lambda-\ell$ and $q\equiv 1-G(1-\kappa_p)$, the share
of detections that proceed.

\begin{lemma}[Credit and overlap]\label{lem:credit}
Under \cref{ass:phi,ass:credit}, expected credit per project is
\begin{equation}\label{eq:credit}
C(R,P)=\exp\!\Big\{-\theta P\big[\lambda-\Delta\,\varphi(R/P)\big]\Big\}.
\end{equation}
Covering papers that the researcher does not pivot away from (\emph{flagged
overlap}) arrive at rate $\theta P\,[1-(1-q)\varphi]$. A share
\begin{equation}\label{eq:aware}
A=\frac{q\,\varphi}{1-(1-q)\,\varphi}
\end{equation}
of them is \emph{aware}: the later paper cites the earlier one. Unaware
duplicates keep a share $1-\lambda$ of credit. Aware follow-ons keep
$\E[\delta\mid\delta>1-\kappa_p]$, which can exceed one.
\end{lemma}

The lemma follows from Poisson thinning. Detected and undetected covering
papers are independent Poisson counts with means $\theta P\varphi$ and
$\theta P(1-\varphi)$, and $\E[t^K]=e^{-\mu(1-t)}$ for $K\sim$ Poisson$(\mu)$.
Three features matter below. Only $\theta\lambda$ and $\Delta/\lambda$ enter
payoffs. Aware follow-ons are selected on high $\delta$, which is why they can
out-earn comparable papers. And the measured overlap flag is a fixed share of
covering events. Its level is a normalisation; its changes are not.

\subsection{The researcher's problem}

A researcher in a topic of density $P$ earns credit $a\,x\,Q(R,P)\,C(R,P)$ per
period. Writing $x$ projects is worth $x$, and each is worth $QC$. Dividing
by the common factor $a$,
\begin{equation}\label{eq:g}
g(r,P)\equiv(1-r)\,Q(br,P)\,C(br,P),\qquad v(P)\equiv\max_{r\in[0,1)}g(r,P).
\end{equation}
The \emph{payoff curve} $v$ is the central object of the paper. Writing speed $a$
scales every payoff and so drops out of $v$. The researcher takes $P$ as given:
she ignores both that her papers teach her colleagues and that they add to
what her colleagues must read.

\begin{remark}[Topic size and what density measures]\label{rem:breadth}
A topic has a fixed set of projects: the coverage rate $\theta$ is a property
of the topic and does not shrink when the topic attracts more people. This is
a normalisation, not a restriction, once density is measured locally. Let a
field of breadth $B$ consist of $B$ cells of unit breadth, each receiving a
flow $P/B$, and let researchers read and be covered within their own cell, so
that the coverage rate per paper is $\theta_0/B$. The field is then $B$
independent copies of a unit topic with density $P/B$, and every result below
holds with $P$ replaced by $P/B$. The model's topic is therefore a
neighbourhood of fixed size in idea space, and $P$ is papers per unit of idea
space: the quantity that nearest-neighbour similarity measures. If papers fall
at rate $P/B$ in a $d$-dimensional space, the distance to the $k$-th nearest
earlier paper scales as $(kB/P)^{1/d}$. ``Opening a topic'' means moving into
an empty cell, and the question of whether topics broaden as they grow becomes
the question of how the volume elasticity varies with the similarity radius,
which \cref{sec:facts} measures.
\end{remark}

\subsection{The frontier, careers and catch-up}

\paragraph{Frontier.} Instead of joining an open topic, a researcher may explore
new ground. An explorer opens a viable topic with probability $\pi(M)$, where $M$
is the measure of topics already open. If she succeeds, her output is worth
$f_0$ per unit of writing speed: the founding papers of a field. The frontier
pays $a\,F(M)$ in expectation, with $F(M)=f_0\,\pi(M)$. We assume $\pi$ is
continuous, positive and nonincreasing, with elasticity
$\eta\equiv -M\pi'(M)/\pi(M)\ge 0$: as more directions are taken, a fruitful
new one is harder to find \citep{kortum1997,bloometal2020}. Opening $M$ topics
takes $\Phi(M)=\int_0^M \dd m/\pi(m)$ researchers. We assume $F(M)<\vhat$, the
peak of the payoff curve defined below, so open topics are worth joining. A
newly opened topic attracts followers only if a thin topic beats the frontier,
$F<v(P_0)$, where $P_0$ is the flow of a newborn topic.

\paragraph{Careers and catch-up.} Researchers leave their topic at rate $\rho$
(retirement, a change of field) and are replaced by newcomers who choose a
topic. The literature a newcomer must master is a stock $S$ that accumulates
the flow and loses relevance at rate $\delta$: $\dot S=P-\delta S$, so
$S=P/\delta$ in a stationary topic. A newcomer must read it at speed $b_c$
before contributing. This is the \emph{catch-up} cost: a topic-level version
of the burden of knowledge of \citet{jones2009}. Explorers face no stock.

\begin{lemma}[Career value]\label{lem:career}
With exponential careers, a newcomer who joins a stationary topic of density
$P$ and then earns $a\,v(P)$ per period obtains, relative to an explorer, the
expected career value $w(P)/F(M)$, where
\begin{equation}\label{eq:w}
w(P)\equiv v(P)\,e^{-\kappa P},\qquad \kappa\equiv\frac{\rho}{\delta\, b_c}.
\end{equation}
In a stationary topic a share $1-e^{-\kappa P}$ of members are catching up,
so $w(P)$ is also output per member and the topic has
$n(P)=P\,e^{\kappa P}/(a\,x^\ast(P))$ members, where $x^*(P)$ is the optimal
writing share.
\end{lemma}

Catch-up takes $T=S/b_c=P/(\delta b_c)$ periods, and a career survives it with
probability $e^{-\rho T}=e^{-\kappa P}$. The frictionless benchmark is
$\kappa=0$ ($b_c\to\infty$). Incumbents earn $a\,v(P)$ per period;
newcomers earn $a\,w(P)$ on average over their careers.

\begin{table}[t]
\centering\footnotesize
\caption{Notation}\label{tab:notation}
\begin{tabularx}{\textwidth}{@{}>{\raggedright\arraybackslash}p{2.3cm}X@{\hspace{1.2em}}>{\raggedright\arraybackslash}p{2.1cm}X@{}}
\toprule
$a$, $b$ & papers written / read per period at full time & $P$ & topic density (flow of papers per period) \\
$x$, $r$ & writing and reading shares, $r=1-x$ & $R=br$, $c=R/P$ & papers read; reading coverage \\
$Q(R,P)$ & quality gain from what is read & $\varphi(c)$ & probability a covering paper is caught \\
$\theta$ & coverage rate per earlier paper & $\lambda$, $\ell$ & credit lost per unread / caught covering paper \\
$g(r,P)$, $v(P)$ & payoff per unit of $a$; its maximum & $\Phat$, $\vhat$ & peak of $v$ and its height \\
$F(M)$, $\pi(M)$ & frontier payoff; explorer's success probability & $\eta$ & elasticity of $\pi$ in topics opened $M$ \\
$\rho$, $\delta$ & career turnover; decay of the stock's relevance & $S$, $b_c$ & stock of literature; catch-up reading speed \\
$\kappa=\rho/(\delta b_c)$ & catch-up cost per unit of density & $w(P)$ & newcomer's career payoff, $v(P)e^{-\kappa P}$ \\
$P^*$, $P^m$ & equilibrium (mature) density & $\mathcal E$ & elasticity of density with respect to $N$ \\
$g$, $g^*$ & topic diffusion rate; overshooting threshold & $\lambda_J$ & stepping-on-toes parameter of \citet{jones1995} \\
$\lambda_s$, $\tau$ & social loss per unread covering paper; transfer share & $U(P)$ & entrant's expected payoff along the consolidation path \\
\bottomrule
\end{tabularx}
\end{table}

\subsection{Equilibrium}

\begin{definition}[Topic equilibrium]\label{def:eq}
A stationary topic equilibrium is a density $P^*$ for every open topic,
a writing share $x^*$ and a measure $M^*$ of open topics such that:
(i) $x^*=x^*(P^*)$ solves \eqref{eq:g};
(ii) newcomers are indifferent between joining any open topic and exploring:
$w(P^*)=F(M^*)$;
(iii) the allocation is stable under entry: $w'(P^*)<0$;
(iv) researchers add up: $\Phi(M^*)+M^*n(P^*)=N$.
\end{definition}

Condition (iii) selects allocations that survive the adjustment
$\dot n_j\propto w(P_j)-F(M)$. At a crossing where $w$ rises, one more entrant
raises the payoff of entering and draws in more. Topics in equilibrium need
not have the same density a priori, but (ii) and (iii) leave a single
admissible density, so they do. We state results in the \emph{large-topic
limit} $\Phi(M)/N\to 0$. Here topics are large relative to the researcher time
needed to open one, as they are in the data, where pioneers are a small share
of authors. \Cref{app:proofs} gives the general expressions.
